\chapter{Einleitung}
Die fortschreitende Digitalisierung und die zunehmende globale Vernetzung von Unternehmen führen zu einer stetig wachsenden Bedeutung effizienter Kommunikations- und Informationsprozesse. Insbesondere im modernen Büroumfeld stellen Meetings ein zentrales Element der Zusammenarbeit dar. Sie dienen der Abstimmung, Entscheidungsfindung und Koordination von Arbeitsabläufen. Gleichzeitig steigt jedoch mit der Anzahl und Dauer von Besprechungen auch der organisatorische Aufwand, der mit deren Vor- und Nachbereitung verbunden ist.

Trotz der Verbreitung moderner Kommunikationsplattformen wie Chat-Systemen oder Videokonferenzlösungen bleibt die strukturierte Dokumentation von Meetings eine Herausforderung. Häufig fehlen Zeit und Ressourcen, um besprochene Inhalte vollständig und nachvollziehbar festzuhalten. In der Praxis führt dies oftmals dazu, dass Informationen verloren gehen, Entscheidungen nicht eindeutig dokumentiert werden oder Themen mehrfach diskutiert werden müssen. Insbesondere in Arbeitsumgebungen mit hoher Taktung an Meetings entsteht dadurch eine erhebliche Informationsdichte, die von den Beteiligten kaum effizient verarbeitet werden kann.

Ein möglicher Ansatz zur Bewältigung dieser Problematik liegt im Einsatz KI-gestützter Systeme zur automatisierten Verarbeitung gesprochener Sprache. Moderne Verfahren der automatischen Spracherkennung und Textverarbeitung ermöglichen es, gesprochene Inhalte in Textform zu überführen, zu strukturieren und zusammenzufassen. In der Praxis zeigen sich jedoch weiterhin Herausforderungen hinsichtlich Genauigkeit, zeitlicher Verfügbarkeit der Ergebnisse sowie der sinnvollen Aufbereitung der Informationen. Insbesondere nachträglich erzeugte Zusammenfassungen sind fehleranfällig oder stehen nicht unmittelbar zur Verfügung, wodurch ihr Nutzen im Arbeitsalltag eingeschränkt wird.

Vor diesem Hintergrund befasst sich diese Masterarbeit mit der Entwicklung eines KI-basierten Meetingassistenten zur automatisierten Erfassung, Aufbereitung und Darstellung gesprochener Inhalte. Ziel ist es, ein System zu konzipieren und umzusetzen, das Meetings in Echtzeit analysiert, relevante Inhalte strukturiert aufbereitet und dadurch die Nachbereitung von Besprechungen deutlich erleichtert. Der Fokus liegt dabei auf der praktischen Umsetzbarkeit sowie der Integration moderner Verfahren aus den Bereichen Sprachverarbeitung und Automatisierungstechnik.

Die Arbeit ist im Studiengang Elektro- und Informationstechnik mit der Vertiefung Automatisierung und Robotik angesiedelt und verbindet theoretische Konzepte der künstlichen Intelligenz mit einer anwendungsnahen technischen Umsetzung. Ziel ist es, einen Beitrag zur Effizienzsteigerung wissensbasierter Arbeitsprozesse zu leisten und das Potenzial KI-gestützter Assistenzsysteme im beruflichen Kontext aufzuzeigen.

\section{Notisent}
Notisent GmbH ist ein 2025 gegründetes Technologie-Start-up mit Sitz in Göttingen, das sich auf die Entwicklung KI-gestützte Arbeitsassistenten für wissensintensive Arbeitsprozesse spezialisiert hat. 
Ziel des Unternehmens ist es, die zunehmende Informationsdichte in digitalen Arbeitsumgebungen durch automatisierte Analyse- und Strukturierungsmechanismen beherrschbar zu machen und dadurch die Effizienz sowie die Qualität wissensbasierter Tätigkeiten zu erhöhen. 
Das Unternehmen entwickelt einen intelligenten digitalen Assistenten, der es Mitarbeitenden und Organisationen ermöglicht, den täglichen Informations- und Kommunikationsfluss effizienter zu bewältigen. Im Zentrum der Lösung steht eine KI-Plattform, die E-Mail, Dokumente, Notizen, Kalenderinformationen und andere Datenquellen analysiert, verknüpft und priorisiert, um arbeitsrelevante Inhalte und Aufgaben gezielt bereitzustellen. 

\section{Gegenstand und Ziel}
Ziel der Arbeit ist die Entwicklung eines Meetingassistenten-Prototypen, welcher in Echtzeit in der Lage ist gesprochene Inhalte in einem Online-Meeting zu analysieren und anschließend mithilfe von KI für den Benutzer grafisch aufzubereiten. Dazu sollen im ersten Kapitel die notwendigen Grundlagen geschaffen werden für Technisches Verständnis. Im zweiten Kapitel Methodik soll dem Leser die herangehensweise an das Projekt vermittelt werden und anschließend im darauffolgenden Kapitel die Umsetzung des Prototypen. 

Der Meetingassistent wird im Webframework Django in der Programmiersprache Python entwickelt. Die Grundlage des Meetingassistenten bildet ein Large Language Modell, welches auf Servern des KI-Servicezentrums für sensible und kritische Infrastruktur in Göttingen gehostet wird. 

Insgesamt liegt der Fokus auf der Entwicklung eines ersten funktionsfähigen Prototypen, welcher in die Notisent Umgebung eingebunden werden kann. In Abbildung X ist eine Übersicht über alle Komponenten dargestellt. 

Grafik wie grundsätzlich funktionieren soll 

Anhand der folgenden Forschungsfragen wird präzisiert, welche Ergebnisse diese Arbeit leifert. 
Forschungsfragen:
\begin{enumerate}
    \item\textbf{Wie kann ein Live-Meetingassistent entwickelt werden?} 

Diese Frage untersucht wie so ein Live-Meetingassistent vom Grundkonzept her aufgebaut werden muss um funktioniern zu können. 

    \item\textbf{Was für Anforderungen muss so ein Meetingassistent erfüllen?}

In diesem Zusammenhang soll untersucht werden, welche Anforderungen ein Meetingassistent erfüllen muss um einen Mehrwert für ein Meeting zu bieten. 

    \item\textbf{Wie ist die Qualität im vergleich zu Komplett zusammengefassten Meetings?} 

Es wird erforscht, wie sich ein Live-Meetingassistent im vergleich zu bereits bestehenden Tools zum zusammenfassen 

    \item\textbf{Kann die Zusammenfassung in annehmbarer Zeit erstellt werden?} 
    
\end{enumerate}

Diese Forschungsfragen bilden die Grundlage für die Entwicklung eines leistungsfähigen, flexiblen und nutzerfreundlichen KI-Meetingassistenten, welcher sich in die Notisent Umgebung integrieren lässt und 
 

 
 


